\documentclass[10pt,a4paper,oneside]{article}

\usepackage{standalone}
\usepackage[utf8]{inputenc}
\usepackage{amsmath}
\usepackage{amsfonts}
\usepackage{amssymb}
\usepackage{graphicx}
\usepackage{tocloft}
\usepackage{hyperref}
\usepackage{float}
\usepackage{darkmode}

\enabledarkmode

\author{
  Vérin Clément \\ \texttt{ISIMA, ZZ2 F1}
  \and
  Scrimieri Daniele \\ \texttt{BradFord University}
}
\title{SoftWare System Design and Testing \\ }

\setlength{\parindent}{0pt}

\hypersetup
{
  colorlinks=true,
  linkcolor=black,
  urlcolor=cyan,
  linkcolor=white
}

\begin{document}

\fbox
{
  \parbox{\textwidth}
  {
    \centering
    \maketitle
  }
}

\vspace{1cm}

\begin{figure}[H]
	\vspace{0cm}
    \centering
    \includegraphics[width=0.8\textwidth]{src/lateX/Captures/Title.png}
    \caption{MFA is a commonly used and secured way to identify yourself}
\end{figure}

\begin{center}
\href{https://bradford.instructure.com/courses/21719/assignments/114047} {Deposit Link}
\end{center}

\renewcommand{\contentsname}{Summary}
\renewcommand{\cftsecleader}{\cftdotfill{\cftdotsep}}
\renewcommand{\cftdotsep}{2}

\newpage

\tableofcontents

\section{Design}
The first work's goal were about extending the previously implemented MFA, Multi-Factor-Authentification situation with an interactive GUI Interface.\\
I recovered my last CourseWork 1 code and added 6 new classes entirely dedicated to GUI.\\

Basically, these GUI classes are arranged into two types regarding what they are related with: Authentification Method GUI or Platform (or OS) GUI.\\

The \textbf{Method\_GUI} and \textbf{Platform\_GUI} are two abstract classes providing the basics a GUI needs.\\
Actually, I thought using abstract classes were the best because Authentification Methods and OS pretty have the same behaviour, so I managed to type a maximum of method in these abstract classes so actual sub-classes only have a few to Override.\\

For example, what's make the difference between OS are only their name, and the type of required Identifcation feature for Authentification method (so Mail or Phone number).\\
I assume this strategy allows a better flexibility, make the code opened to easy extension (\textit{Open Closed Pinciple}) and minimize code duplication.\\

Besides, almost every class depends on interface or abstract classes, which follows the \textit{Interface Segregation Principle}.\\
However, some classes still have more than one role or purpose which do not fully respects the \textit{Interface Responsibility Principle} from SOLID.\\

\newpage
Here is the resulting UML Diagram, which allows a better visual understanding :

\begin{figure}[H]
	\vspace{0cm}
    \centering
    \includegraphics[width=1.2\textwidth]{src/uml/MFA\_GUI.png}
    \caption{UML Diagram for MFA GUI extension}
\end{figure}

(if you got some troubles reading the diagram, you can find it in \textit{src/uml}).

Of course, Iv'e adapted the original code by getting rid of Scanner(System.in) waiter, and I opted for CountDownLatch which allows us to stop a code and wait for the user to enter something from the GUI.\\
I used \textit{Java Swing} but I still have some trouble with it: sometimes, when I update my evoluting \textit{JPanel}, new features doesn't shows up.\\
The only solution for that is too set the visibility of the window to false, then true, and features will appear.

\section{Implementation}
You can check the \textit{makefile} in order to know wich possibility you've got.\\

But you can just type \textit{make run\_GUI\_MFA} \textit{or java -cp target/classes MFA\_Terminal.Main}, and \textit{make run\_Terminal\_MFA} in order to run the original CourseWork1.\\
You can also type \textit{make test} in order to launch the test series.\\

You can find some documentation in the \textit{target/generated\_ressources/documentation} folder.\\

\newpage
Here is an example of what you could see running the GUI :

\begin{figure}[H]
	\vspace{0cm}
    \centering
    \includegraphics[width=1.2\textwidth]{src/lateX/Captures/C1.png}
    \caption{GUI Running with Windows OS and Mail code authentification method (Sucess)}
\end{figure}

\begin{figure}[H]
	\vspace{0cm}
    \centering
    \includegraphics[width=1.2\textwidth]{src/lateX/Captures/C2.png}
    \caption{GUI Running with Linux OS and SMS code (or call resonding) authentification method (Failed)}
\end{figure}

\begin{figure}[H]
	\vspace{0cm}
    \centering
    \includegraphics[width=1.2\textwidth]{src/lateX/Captures/C3.png}
    \caption{All the possible accounts you can use}
\end{figure}

\section{Questions}

\subsection{A}
Factory design is a pretty useful pattern which allows us to create classes or instances regardless their "types" or "specificities", bringing more flexibility and adaptability to our code.\\

\textbf{Pros}:
- The created class is himself Encapsulated: Client doesn’t need to know wihch exact class has been made or how it’s constructed.\\
- "Easy Decoupling": differentiating class creators provide less tight coupling or strong dependencies.\\
- Flexibility: Makes it easy to introduce new subclasses without changing client code (Open/Closed Principle from SOLID).
- Ordering: all calsse's basic code is "centralized" in one place, enhancing flexibility.\\

\textbf{Cons}:
- Complexity: This pattern could introduce a lot of interface or abstract classes that may feel unnecessary for easy objects.\\
- Override: Factory pattern could induce a lot of Overheading wich can make the work harder for both machine and developper.\\
- Hard debugging: factory pattern can oftenly hide concrete classes, which can make Instantiation even harder sometimes.\\

\textbf{When to use Factories}:
The best situation to use factory pattern seems when the exact type of object to create is unknown until or during runtime.\\
Or when a lot of versions of the same concept could exist (Example: Dishes in restaurant).\\

\subsection{B}
Dependencies represents the external components that a class or system relies on to operate (could be databases, APIs, utility classes, etc.)\\

dependencies are important in software architecture because:\\
- They define how modules communicate, interact and classify themselves.\\
- Bad managed dependencies can lead to tight coupling, obligation to making changes, hard testing and re-usability.\\
- Well managed dependencies improves modularity, scalability, and testability.\\

We shall depend on abstract interfaces and not concrete implementations because SOLID: \textit{Dependency Inversion Principle}.\\
"Depend on abstractions. Do not depend on concretions."

Using abstract interfaces provides an better testability by knowing about dependencies during unit tests.\\
Abstract interfaces as well enable us to use more polymorphism and composition.\\

\section{Unit Testing}
You can find test java source in \textit{src/test/java/SRC\_MFA\_TERMINAL\_TEST\_DIR} (which are the same than before, I did not figure how to properly test GUI features ...)\\
You can as well find everything related to Jacoco Coverage in \textit{target/generated\_ressources/jacoco/jacoco-report}.\\

You can also type \textit{make test} in order to launch the test series, or \textit{make coverage} to re-generate the Jacoco Coverage.\\

\begin{figure}[H]
	\vspace{0cm}
    \centering
    \includegraphics[width=0.8\textwidth]{src/lateX/Captures/C4.png}
    \caption{Jacoco Report}
\end{figure}

(The coverage is really really bad actually, because of untested GUI)

\section{Mutation Testing}

\subsection{A}
For mutation Testing, I chose \textbf{PITEST}.\\

In order ot use Pitest, I have downloaded 3 \textit{.jar file} containing every needed: \textit{pitest-1.15.2.jar, pitest-command-line-1.15.2.jar and pitest-entry-1.15.2.jar}.\\
And only then, I can start launching Mutation programm by typing  \textit{make mutation} (this is a pretty big command line actually).

The only problem is: Pitest keep saying "No mutators has been found" ...

\begin{figure}[H]
	\vspace{0cm}
    \centering
    \includegraphics[width=0.8\textwidth]{src/lateX/Captures/C5.png}
    \caption{Pitest " No Mutation Find" exception}
\end{figure}

I tried a lot of solution. I don't think the issue comes from the source test.\\

\subsection{B}
\subsection{C}

\section{Conclusion}
{
  This Coursework definitely improved my understanding of Java Swing, including some of his components and features.\\
  But coding some Java Swing can take a lot of line to deal with, so it encouraged me to use a little bit more than usual Encapsulation Prinicples and Rules which drastically reduced the amount of code.\\

  However, I still struggle to make godd and interesting test, but I am glad I learned about Coverage and Mutation programms which definitely can be useful in the future, even if I didn't manage to use Pitest.\\
}
\end{document}